\chapter{導入}
\label{sec:introduction}

\section{研究の背景}
情報アクセスシステムは，利用者の問いや作業に関連する資料を探すために用いられる．
検索結果を順位付きの一覧として提示する場合もあれば，取得した資料を要約して回答を作る場合もある．
いずれの場合も，何をもって適切な情報と判断するかを明確にしなければ，システムの比較は難しい．
同じ資料であっても，利用者が求める情報の範囲や詳しさによって，役立ち方が異なるためである．

本サンプルでは，検索結果の並べ替えを題材として，修士論文の構成と記述方法を示す．
比較するのは，最初の検索順位をそのまま用いる基準システムと，候補文書に別の得点を与えて
順位を変更するシステムである．ただし，ここに記す手順は説明用の設計例であり，実装や実験を
実施したという意味ではない．実際の論文では，計画と実施済みの内容を区別して記述する必要がある．

\section{研究課題}
検討する課題は二つある．第一に，追加の並べ替えによって上位の検索結果が改善するかを調べる．
第二に，改善が特定の検索要求に偏っていないかを確認する．平均値が上昇したとしても，
すべての検索要求で同じ傾向が得られるとは限らない．一部の要求で大きな改善があり，
別の要求では悪化するという組合せでも，平均値は上昇し得る．

この違いを確認するため，比較では検索要求ごとの得点を保存する．文書集合，検索要求，
適合性判定を両方のシステムで共通にし，変更点を順位付けの処理に限定する．
これにより比較の条件は整理しやすくなるが，評価資料の不足や要求の選び方による不確実性が
なくなるわけではない．結果を解釈する際には，そのような制約も合わせて検討する．

% Demonstration-only break: expose a continuation page.
\newpage
\section{対象範囲と前提}
\label{sec:scope}
本例では，固定された文書集合に対して検索を行い，あらかじめ用意された判定を使って結果を評価する．
実験の途中で文書を追加したり，片方のシステムだけ異なる判定を使用したりしない．
同一条件での比較を行うためには，使用した資料の名称だけでなく，版や更新日も記録することが重要である．
同じ名称の資料であっても，内容が更新されていれば結果が一致しない可能性がある．

評価では，平均的な性能と個々の検索要求の挙動を分けて確認する．平均値は全体を簡潔に示せる一方，
得点の分布を直接表すものではない．個別の結果を保存しておけば，平均値が変化した理由を追跡し，
入力の欠落や識別子の不一致なども確認できる．このような記録は，結果が期待に沿わなかった場合にも必要となる．

\subsection{比較の単位}
比較の単位は一つの検索要求とする．各要求に対して両方のシステムの得点を求め，
その差を記録する．一方の結果が欠けている場合には，単にその要求を取り除くのではなく，
欠落した理由と評価上の取扱いを明記する．除外によって比較対象が変われば，得られた平均値の意味も変わるためである．

\subsection{主張の範囲}
固定された資料を用いた評価から，実際の利用場面での使いやすさを直接結論付けることはできない．
画面の構成，応答時間，利用者の知識など，別の要因も情報の利用に影響する．
本例で扱うのは順位付けの比較手順であり，人を対象とした実験の結果ではない．
結論では，どの条件の下で得られた知見なのかを改めて示す必要がある．

% A second continuation page exposes the other odd/even header position.
\newpage
\section{論文の構成}
第\ref{sec:related_work}章では，関連研究を整理する際の観点を示す．
第\ref{sec:experiments}章では，比較条件，集計方法，表や図の記載例を説明する．
第\ref{sec:approach}章では，評価の問題点と，それに対応する手順の改善案を整理する．
付録\ref{app:details}には，再現に必要な情報と，結果を点検するための記録例を示す．
参照先の番号は自動的に更新されるため，章を追加した場合にも番号を手入力する必要はない．

各章の役割は研究内容に応じて調整してよい．例えば，実装上の工夫が主要な貢献であれば，
システムの構成を説明する独立した章を設ける方法がある．複数の実験が異なる研究課題に対応する場合には，
課題ごとに方法と結果をまとめる構成も考えられる．テンプレートの章立てをそのまま維持することよりも，
問いと証拠の関係を読み手が追えることが重要である．

\section{書式の確認方法}
このページでは，ページ上部に章または節の見出し，下部中央にページ番号が表示される．
偶数ページには章の見出しを左側に，奇数ページには節の見出しを右側に表示する．
一方，章の最初のページでは，上部の見出しと横線を表示せず，ページ番号のみを残す．
この章を三ページ構成とすることで，両方の通常ページの表示を確認できるようにしている．

次の章は，この章の直後のページから始まる．開始ページが偶数であっても空白ページを挿入しない．
サンプル内の明示的な改ページは通常ページの表示を確認するために設けたものである．
実際の執筆時にはこれらを取り除き，本文を自然に流してよい．章の開始時の改ページは，
章を指定する命令によって自動的に行われる．
